%! AP Statistics — Unit 4, Lesson 4.11 — Lesson Plan
\documentclass[10pt]{article}
\usepackage{apstats-article}
\usepackage{apstats-boxes}
\usepackage{graphicx}
\graphicspath{{images/}}

\newcommand{\UnitNumberName}{Unit 4: Probability, Random Variables, and Probability Distributions \quad}
\newcommand{\LessonNumberName}{Lesson 4.11: Parameters for a Binomial Distribution}

\begin{document}
\begin{center}
    {\Large\bfseries \CourseName: \SchoolYear \\
    {\small\bfseries \UnitNumberName \LessonNumberName}}
\end{center}

\begin{tcolorbox}[colback=sky, colframe=navy, boxrule=0.9pt, arc=2mm,
    left=2mm, right=2mm, top=1.5mm, bottom=1.5mm]
    \textbf{Primary Objective:} Students will calculate and interpret the mean and
    standard deviation of a binomial distribution using the formulas $\mu_X = np$
    and $\sigma_X = \sqrt{np(1-p)}$, and compute cumulative binomial probabilities
    using \texttt{binomcdf} (Big Idea: UNC-3).
\end{tcolorbox}

\renewcommand{\arraystretch}{1.6}
\begin{skillbox}[Priority Ideas \& Skills]{goldbox}
\begin{minipage}[t]{0.48\linewidth}
    \textbf{AP Skills 3.B and 4.B}
    \begin{itemize}
        \item Calculate $\mu_X = np$ and $\sigma_X = \sqrt{np(1-p)}$ for
              $X \sim B(n, p)$ (UNC-3.B, Skill 3.B).
        \item Interpret $\mu_X$ in context: ``In many repetitions of this
              binomial experiment, the average number of successes would be
              $np$.'' (UNC-3.C, Skill 4.B).
        \item Compute cumulative probabilities $P(X \le k)$ with
              \texttt{binomcdf(n, p, k)} and $P(X \ge k)$ via complement.
    \end{itemize}
\end{minipage}
\hfill
\begin{minipage}[t]{0.48\linewidth}
    \textbf{Key Understandings}
    \begin{itemize}
        \item $\mu_X = np$ and $\sigma_X = \sqrt{np(1-p)}$ are derived from the
              general formulas for a discrete RV, applied to the binomial.
        \item \texttt{binomcdf(n, p, k)} $= P(X \le k)$ (adds up from $0$ to $k$).
        \item $P(X \ge k) = 1 - P(X \le k-1) = 1 - \texttt{binomcdf(n, p, k-1)}$.
        \item On the AP exam, always show the formula setup for $\mu_X$ and
              $\sigma_X$ before plugging in numbers.
    \end{itemize}
\end{minipage}
\end{skillbox}

\begin{skillbox}[{Vocabulary, Concepts, \& Theorems}]{greenbox}
\begin{minipage}[t]{0.48\linewidth}
    \begin{itemize}
        \item \textbf{Mean of a binomial} --- $\mu_X = np$
        \item \textbf{Standard deviation of a binomial} ---
              $\sigma_X = \sqrt{np(1-p)}$
        \item \textbf{Cumulative binomial probability} ---
              $P(X \le k) = \sum_{j=0}^{k} \binom{n}{j}p^j(1-p)^{n-j}$
    \end{itemize}
\end{minipage}
\hfill
\begin{minipage}[t]{0.48\linewidth}
    \begin{itemize}
        \item \textbf{binomcdf(n, p, k)} --- calculator command for $P(X \le k)$
        \item \textbf{Complement rule for binomcdf} ---
              $P(X \ge k) = 1 - \texttt{binomcdf(n, p, k-1)}$
        \item \textbf{Interpretation of $\mu_X$} --- long-run average number of
              successes per experiment
    \end{itemize}
\end{minipage}
\end{skillbox}

\begin{skillbox}[Activate Prior Knowledge \& Spiral Review]{sky}
\begin{minipage}[t]{0.48\linewidth}
    \textbf{Prior Knowledge Check (3--4 min)}
    \begin{itemize}
        \item Review: How do we interpret the mean $\mu_X$ of any discrete random
              variable? (Lesson 4.6)
        \item Connect: Last lesson we computed $P(X = k)$; today we summarize the
              whole distribution with just two numbers: $\mu_X$ and $\sigma_X$.
        \item Prompt: \textit{``If a binomial experiment has $n = 20$ trials and
              $p = 0.5$, what would you expect the average number of successes
              to be? How did you get that?''}
    \end{itemize}
\end{minipage}
\hfill
\begin{minipage}[t]{0.48\linewidth}
    \textbf{Connections Forward}
    \begin{itemize}
        \item Geometric mean $\mu_X = 1/p$ $\to$ Lesson 4.12
        \item Large-sample normal approximation to binomial $\to$ Unit 5
        \item Sampling distributions: $\mu_{\hat{p}} = p$,
              $\sigma_{\hat{p}} = \sqrt{p(1-p)/n}$ mirror these formulas $\to$
              Unit 5
    \end{itemize}
\end{minipage}
\end{skillbox}

\begin{skillbox}[Hook]{sky}
    \textbf{Hook (4--5 min)}\\
    Display: \textit{``A professional archer hits the target 90\% of the time.
    She shoots 50 arrows. About how many do you expect her to hit?
    How spread out would her totals be across many competitions?''}\\[4pt]
    Most students will guess $45$ (correct!). Ask them how they reasoned that.
    Draw out the idea of multiplying $n \times p$.
    Then ask about spread --- introduce the idea that $\sigma_X$ quantifies
    variability around that average.
    Transition: \textit{``Today we derive and practice exact formulas for both
    the center and the spread of a binomial distribution.''}
\end{skillbox}

\begin{skillbox}[Lesson]{sky}
\begin{minipage}[t]{0.48\linewidth}
    \textbf{Part 1 --- Mean and SD Formulas (8 min)}
    \begin{enumerate}
        \item State $\mu_X = np$ and $\sigma_X = \sqrt{np(1-p)}$; briefly show
              these come from applying the general discrete RV formulas to a
              binomial (optional derivation).
        \item Model with the archer: $\mu_X = 50(0.9) = 45$;
              $\sigma_X = \sqrt{50(0.9)(0.1)} \approx 2.12$.
        \item Interpret $\mu_X$: \textit{``If she competes in many 50-shot
              competitions, she will average 45 hits per competition.''}
        \item Interpret $\sigma_X$: typical deviation from the mean.
    \end{enumerate}
\end{minipage}
\hfill
\begin{minipage}[t]{0.48\linewidth}
    \textbf{Part 2 --- Cumulative Probabilities (8--10 min)}
    \begin{enumerate}
        \item Define $P(X \le k)$ and explain why it's impractical to add
              individual probabilities by hand for large $n$.
        \item Calculator: \texttt{binomcdf(n, p, k)} gives the cumulative sum.
        \item Model: $P(X \ge 48) = 1 - \texttt{binomcdf(50, 0.9, 47)}$.
        \item AP exam note: show each step ---
              write ``$1 - P(X \le 47)$'' before invoking \texttt{binomcdf}.
        \item Boundary attention: $P(X > k) = 1 - \texttt{binomcdf(n, p, k)}$
              vs.\ $P(X \ge k) = 1 - \texttt{binomcdf(n, p, k-1)}$.
    \end{enumerate}
\end{minipage}
\end{skillbox}

\begin{skillbox}[Active Monitoring]{redbox}
\begin{minipage}[t]{0.48\linewidth}
    \textbf{Circulate \& Check For\ldots}
    \begin{itemize}
        \item Forgetting the square root when computing $\sigma_X$.
        \item Off-by-one errors: $P(X \ge k)$ vs.\ $P(X > k)$ --- require
              students to write the inequality before using \texttt{binomcdf}.
        \item Interpretation errors: $\mu_X$ describes long-run average,
              not a guaranteed outcome for any single experiment.
    \end{itemize}
\end{minipage}
\hfill
\begin{minipage}[t]{0.48\linewidth}
    \textbf{Cold-Call Prompts}
    \begin{itemize}
        \item ``Interpret $\mu_X = 45$ in a sentence a non-statistics student
              could understand.''
        \item ``Why is $P(X \ge 48)$ not the same as
              $1 - \texttt{binomcdf(50, 0.9, 48)}$?''
        \item ``If $p$ were 0.5 instead of 0.9, would $\sigma_X$ be larger
              or smaller? Why?''
    \end{itemize}
\end{minipage}
\end{skillbox}

\begin{skillbox}[Group Work \& Differentiation]{redbox}
\begin{minipage}[t]{0.48\linewidth}
    \textbf{Three-Tier Activity (10--12 min)}
    \begin{itemize}
        \item \textbf{Tier R}: Given $n$ and $p$, compute $\mu_X$ and $\sigma_X$
              and fill in a complete sentence interpretation of $\mu_X$.
        \item \textbf{Tier A}: Compute $\mu_X$, $\sigma_X$, and a cumulative
              probability from a contextual scenario.
        \item \textbf{Tier E}: Compare two binomial distributions by computing
              their means and SDs; determine which is more variable and explain why.
    \end{itemize}
\end{minipage}
\hfill
\begin{minipage}[t]{0.48\linewidth}
    \textbf{Differentiation}
    \begin{itemize}
        \item \textit{Support}: Formula strip with $\mu_X = np$ and
              $\sigma_X = \sqrt{np(1-p)}$; step-by-step model computation.
        \item \textit{Extension}: Explore how $\sigma_X$ varies as $p$ changes
              from 0 to 1 for a fixed $n$; find the value of $p$ that maximizes
              $\sigma_X$.
        \item \textit{ELL}: Sentence frames for interpretation prompts.
    \end{itemize}
\end{minipage}
\end{skillbox}

\begin{skillbox}[Individual Work \& Assessment]{redbox}
\begin{minipage}[t]{0.48\linewidth}
    \textbf{Exit Ticket (5 min)}\\
    Scenario: \textit{``A manufacturer claims 3\% of its products are defective.
    A store orders 200 items.''} Students compute $\mu_X$, $\sigma_X$,
    $P(X \le 10)$, and interpret $\mu_X$ in context.
\end{minipage}
\hfill
\begin{minipage}[t]{0.48\linewidth}
    \textbf{AP-Style Check}\\
    \textit{``For $X \sim B(40, 0.25)$, which gives $P(X \ge 12)$?''}
    \begin{enumerate}[label=(\Alph*)]
        \item $1 - \texttt{binomcdf(40, 0.25, 12)}$
        \item $\texttt{binomcdf(40, 0.25, 12)}$
        \item $1 - \texttt{binomcdf(40, 0.25, 11)}$
        \item $\texttt{binompdf(40, 0.25, 12)}$
    \end{enumerate}
    Answer: (C)
\end{minipage}
\end{skillbox}

\begin{skillbox}[Reinforcement \& Extension]{goldbox}
\begin{minipage}[t]{0.48\linewidth}
    \textbf{Homework / Practice}
    \begin{itemize}
        \item Four scenarios: compute $\mu_X$, $\sigma_X$, and interpret each
              in context.
        \item Two cumulative probability problems using \texttt{binomcdf}.
        \item AP free-response: full solution showing formula, calculator
              notation, and interpretation.
    \end{itemize}
\end{minipage}
\hfill
\begin{minipage}[t]{0.48\linewidth}
    \textbf{Extension / Preview}
    \begin{itemize}
        \item \textit{Extension}: Show that $\sigma_X$ is maximized when
              $p = 0.5$ for any fixed $n$. Why does this make intuitive sense?
        \item \textit{Preview}: Lesson 4.12 introduces the geometric distribution
              --- a setting where we count trials until the \emph{first} success
              rather than fixing $n$ in advance.
    \end{itemize}
\end{minipage}
\end{skillbox}

\end{document}
